\section{Capítulo~\ref{sec:recognition}. Reconhecimento}

A velocidade do reconhecimento, a construção dos primeiros estágios, a leitura linear da resposta e o aprendizado pela continuidade do tempo foram medidos em humanos e macacos; a redução de toda a cadeia a duas operações e o aprendizado com vídeo foram testados nos modelos do livro, e as explicações das ilusões vão das bem fundamentadas às controversas.

{\footnotesize
\setlength{\tabcolsep}{3pt}\renewcommand{\arraystretch}{1.15}
\begin{longtable}{>{\raggedright}p{0.5cm}>{\raggedright}p{4.0cm}>{\raggedright}p{5.6cm}>{\raggedright}p{2.3cm}>{\raggedright\arraybackslash}p{1.7cm}}
\toprule
N.º & Afirmação & Experimento e o que foi medido & Fonte & Status \\
\midrule\endfirsthead
\toprule
N.º & Afirmação & Experimento e o que foi medido & Fonte & Status \\
\midrule\endhead
1 & A comparação pixel a pixel com um modelo, sob deslocamento, não acerta mais que uma moeda; o par de estágios~\eqref{eq:scstage} reconhece a figura em qualquer lugar & \texttt{models/\allowbreak recognition\_models.py}, ``retina'' de $24$ pontos, $4000$ tentativas: modelo, $0.49$, $0.51$, $0.52$ com fração de pontos invertidos $0$, $0.1$, $0.2$; par $S$ e $C$, $1.00$, $0.86$, $0.70$ & dedução no capítulo & modelo \\
2 & Um animal numa imagem é reconhecido em $\sim150$\,ms, numa passagem por uma dezena de estágios de $10$--$20$\,ms & Humanos, tarefa ``há um animal ou não'' em fotos mostradas por $20$\,ms: os potenciais evocados das duas respostas divergem após $\sim150$\,ms. Macaco: a latência da primeira resposta cresce de estágio em estágio (V1 primeiro, depois V2 e V4). O número de estágios e ``zero ou um disparo por estágio'' são deduções desses números & \cite{Thorpe1996,Schmolesky1998} & medido \\
3 & O estágio $S$ é um E sobre as partes, o estágio $C$, o máximo sobre as posições (proposição~\ref{prop:conveyor}) & Gato, córtex visual primário: as células simples respondem a uma borda de certa orientação num certo lugar, as complexas, à mesma orientação numa área maior. Registro intracelular de células complexas: a resposta a duas barras, em muitas células, é próxima da maior das duas respostas, em média mais próxima da operação max. Máximo por inibição mútua: proposição~\ref{prop:wta}; divisão pela atividade total: revisão & \cite{HubelWiesel1962,Lampl2004,Carandini2012} & medido \\
4 & Mais acima na cadeia, as combinações são maiores e mais complexas, e a resposta depende cada vez menos da posição & Macaco, simplificação dos estímulos até o ``atributo crítico'': as células de V4 e do córtex inferotemporal posterior respondem a combinações complexas de forma, cor e textura com campo pequeno; no córtex inferotemporal anterior o campo é maior. A alternância de $S$ e $C$ no modelo reproduz esse quadro & \cite{KobatakeTanaka1994,Riesenhuber1999} & medido \\
5 & As redes convolucionais repetem essa alternância e são as que melhor predizem as respostas dos estágios superiores; o objeto é lido linearmente das taxas de um grupo de neurônios & Rede treinada para reconhecer, sem ajuste ao cérebro: a camada superior prediz as respostas dos neurônios do córtex inferotemporal do macaco, as intermediárias, as de V4. Um classificador linear sobre a atividade de $\sim100$ células escolhidas ao acaso, em janelas a partir de $12.5$\,ms, reconhece o objeto e a categoria com posição e tamanho diferentes & \cite{Fukushima1980,Yamins2014,Hung2005} & medido \\
6 & A regra de Hebb com traço~\eqref{eq:traceinvariance} ensina a invariância com vídeo sem rótulos, se o traço não for mais curto que $\sim20$\,ms & A ideia dos atributos lentos e da regra com traço é teoria. \texttt{models/\allowbreak recognition\_models.py}, dois neurônios $C$, $16$ entradas, $10$ rodadas: pausa entre objetos de $0.3$\,s. Com $\tau_{\text{tr}}=5$\,ms, coincidência com a figura de $0.52$ em média, com $10$\,ms, $0.56$; com $20$, $50$ e $200$\,ms, $1.00$ em todas as rodadas (com $30$\,ms, uma rodada em dez erra) & \cite{Foldiak1991,WiskottSejnowski2002} & modelo \\
7 & No cérebro, a invariância é aprendida pela continuidade do tempo & Macaco: o objeto era trocado discretamente durante um salto do olho para um certo lugar do campo; a invariância de posição dos neurônios do córtex inferotemporal mudava de acordo com a troca, de forma significativa já após uma hora. Que essa seja a regra principal do aprendizado visual é hipótese & \cite{LiDiCarlo2008} & medido \\
8 & Movimento: detectores de direção, depois o mesmo par E/OU sobre áreas grandes & Detectores de direção na retina do coelho; na área MT do macaco, todos os $110$ neurônios registrados são seletivos à direção, e a resposta ao movimento é mais forte que em V1 & \cite{Barlow1965,Albright1984} & medido \\
9 & Os estágios superiores enviam para baixo uma predição, e para cima sobe o erro & O modelo explica os efeitos extraclássicos dos campos do córtex visual primário; observações isoladas concordam (revisão); como construção geral da linha de montagem, não está provado & \cite{RaoBallard1999,Keller2018} & hipótese \\
10 & Imagens difíceis exigem realimentação e várias voltas & Macaco: para imagens ``difíceis'', a decisão no córtex inferotemporal se forma $\sim30$\,ms mais tarde; essas respostas tardias são mal preditas por uma rede direta e melhor por redes com realimentação ou mais profundas & \cite{Kar2019} & medido \\
11 & Uma falsificação imperceptível de pixels engana a rede; a visão humana é muito mais robusta, mas não invulnerável & Em redes: uma pequena perturbação escolhida de propósito muda a resposta; o exemplo ``panda $\to$ gibão'' é do segundo trabalho. Em humanos, com apresentação breve, falsificações que se transferem bem entre redes desviam a escolha & \cite{Szegedy2014,Goodfellow2015,Elsayed2018} & medido \\
12 & Ilusões de expectativa: a entrada pouco confiável cede à expectativa; o pálido se move ``mais devagar'', o vestido é visto de formas diferentes (seção~\ref{sec:illusions}) & Grades de menor contraste parecem se mover mais devagar. Enquete com $1401$ pessoas: o vestido é visto em três variantes estáveis, e uma dica sobre a iluminação troca a cor; em $\sim13\,000$ observadores, o que decide é a suposição sobre a iluminação. Um modelo bayesiano com expectativa de movimentos lentos explica várias ilusões de movimento & \cite{Thompson1982,LaferSousa2015,Wallisch2017,Weiss2002,Purves2011} & hipótese \\
13 & Controle de ganho: cachoeira, inclinação e contraste; a grade de Hermann não é inibição dos vizinhos & Os detectores de direção da retina do coelho, com movimento prolongado, reduzem a taxa e, depois da parada, caem abaixo do fundo. A ilusão de inclinação é reproduzida por um modelo com divisão pela atividade dos vizinhos. Mudanças na grade que não alteram a resposta dos neurônios de saída da retina eliminam as manchas: a explicação pela inibição dos vizinhos foi rejeitada & \cite{BarlowHill1963,Schwartz2009,SchillerCarvey2005} & medido \\
14 & Compensação de atrasos: um objeto em movimento parece à frente do clarão & A ilusão foi medida. A psicofísica contradiz tanto o prolongamento do movimento quanto a diferença de latências; foi proposta uma explicação retroativa, pelos eventos $\sim80$\,ms depois do clarão. A disputa não está resolvida & \cite{Nijhawan1994,EaglemanSejnowski2000} & hipótese \\
15 & ``Serpentes'' paradas giram: a diferença de latências~\eqref{eq:latency} é lida pelos detectores de direção & A ilusão funciona também sem cor; pares de elementos vizinhos a geram isoladamente; neurônios do córtex do macaco seletivos à direção dão resposta direcional a esses pares parados. Em humanos, a rotação é disparada por microssacadas e piscadas & \cite{Conway2005,OteroMillan2012} & medido \\
16 & Imagens ambíguas: inibição mútua, fadiga e ruído; as trocas são causadas principalmente pelo ruído & Modelo de grupos de neurônios em competição: nele, as trocas são causadas pelo \emph{ruído}, e sem ruído não existem; ele explica o aumento da frequência das trocas com a força do estímulo. A fadiga tem aqui um papel menor & \cite{MorenoBote2007} & hipótese \\
17 & Parte das ilusões é consequência da tarefa que a máquina aprende & Uma rede treinada para prever o próximo quadro de vídeo ``vê'' rotação em anéis parados e não a vê nas imagens de controle; redes para limpar imagens e aumentar a nitidez caem em ilusões de cor e brilho, como os humanos & \cite{Watanabe2018,GomezVilla2020} & medido \\
\bottomrule
\end{longtable}}
